%% ============================================================
%%  Section: Identification Strategy
%%  Depositor Discipline in Modern Banking (2026-04-07 draft)
%%
%%  Changes from 20260404 version:
%%    - Removed all [@@@ ...] placeholders and addressed each one.
%%    - Toned down "silver bullet" / overselling language throughout.
%%      "instrument" → "natural experiment"; softened "precisely identified"
%%      and "orthogonal" claims.
%%    - Removed all \paragraph{} headers; content flows as prose.
%%    - Fewer em-dashes; replaced with commas or parentheses.
%%    - Staggered-adoption paragraph shortened; removed Callaway/Sun/
%%      de Chaisemartin citations and removed stacked-cohort robustness
%%      claims that are not in the code.
%%    - Phase-in election paragraph kept (mechanism-test role) but shortened.
%%    - Threats to Identification: collapsed to one paragraph; removed
%%      Anticipation paragraph and General Equilibrium paragraph per user
%%      instruction.
%%    - Signal Validation: shortened by ~half; removed \paragraph{} headers
%%      and detailed coefficient recitation; retained key findings.
%%    - Removed the BIB ENTRIES comment block at end of file.
%%
%%  Cross-references:
%%    eq:dayone  (sec:background)
%%    tab:descriptive_sumstats_cecl, tab:cecl_predictors_cross_section
%%    fig:cecl_equity_distribution
%%    fig:es_large_banks_market_cds
%%    sec:results, sec:results:heterogeneity
%% ============================================================

\section{Identification Strategy}
\label{sec:identification}

Our goal is to identify the causal effect of a credit-risk disclosure on
depositor behavior.  The empirical obstacle is familiar: riskier banks
carry fewer uninsured deposits in equilibrium, so cross-sectional comparisons
conflate depositors' information responses with pre-existing sorting and
equilibrium co-determination \citep{covitz2004market, chen2022bank}.
We address this problem by exploiting the CECL Day-One adjustment as a
natural experiment that delivers a well-defined information shock, one whose
timing is set by regulation and whose magnitude reflects historical lending
decisions rather than current deposit market conditions.  

\subsection{Empirical Design}
\label{sec:identification:did}

Our baseline specification is a difference-in-differences estimator that
exploits cross-sectional variation in CECL exposure and the common
regulatory adoption date:
\begin{equation}
  y_{it}
  = \alpha_i + \gamma_t
  + \beta \cdot \mathrm{Post}_{it} \times \mathrm{CECL}_i
  + \mathbf{X}_{it-1}'\boldsymbol{\delta}
  + \varepsilon_{it},
  \label{eq:did}
\end{equation}
where $y_{it}$ is the outcome for bank $i$ in quarter $t$; $\alpha_i$ and
$\gamma_t$ are bank and calendar-quarter fixed effects; $\mathrm{Post}_{it}$
equals one for all quarters at or after bank $i$'s CECL adoption date;
$\mathrm{CECL}_i$ is either the continuous \texttt{CECL Adj} measure or the
binary \texttt{High CECL} indicator, both defined in
Section~\ref{sec:background}; and $\mathbf{X}_{it-1}$ is a vector of lagged
controls (log assets, equity-to-assets ratio, interest expense-to-assets).
Bank fixed effects absorb all time-invariant differences across institutions,
including balance-sheet composition, business model, geographic market, and
charter type, so that $\hat{\beta}$ is identified from within-bank changes
in outcomes over time.  Quarter fixed effects absorb any aggregate
macroeconomic force common to all banks in a given quarter, including the
Federal Reserve's tightening cycle, aggregate deposit normalization, and any
nationwide reaction to the 2023 banking stress.

To evaluate the parallel-trends assumption and trace the dynamics of depositor
responses, we also estimate an event-study specification:
\begin{equation}
  y_{it}
  = \alpha_i + \gamma_t
  + \sum_{\substack{k=-10 \\ k \neq -1}}^{10}
    \beta_k \cdot \mathbf{1}\!\left\{\tau_{it} = k\right\}
    \times \mathrm{CECL}_i
  + \mathbf{X}_{it-1}'\boldsymbol{\delta}
  + \varepsilon_{it},
  \label{eq:eventstudy}
\end{equation}
where $\tau_{it} = t - t_i^{*}$ is event time in quarters relative to bank
$i$'s CECL adoption date, and $k = -1$ is the omitted reference period.
Coefficients $\{\hat{\beta}_k\}_{k < 0}$ test for pre-adoption differential
trends between high- and low-CECL banks; coefficients $\{\hat{\beta}_k\}_{k
\geq 0}$ trace the cumulative post-adoption effect.  A pattern of flat
pre-trends followed by a post-adoption break supports the interpretation that
observed differences are caused by the CECL disclosure rather than
pre-existing divergence.


\subsection{Treatment Exogeneity}
\label{sec:identification:predetermination}

For $\hat{\beta}$ in \eqref{eq:did} to have a causal interpretation,
the CECL treatment must be uncorrelated with potential outcomes absent the
disclosure.  This requires that neither the timing of the shock nor its
magnitude was determined by, or correlated with, the deposit dynamics we
study.

The 2023Q1 adoption date was set by FASB's ASU 2016-13 and enforced by the
federal banking agencies.  Institutions subject to the community-bank rollout
had no ability to accelerate or delay adoption to coincide with favorable
deposit market conditions, so the adoption date is unrelated to any
bank-specific information about contemporaneous deposit flows.

The size of the Day-One adjustment reflects the cumulative composition of the
loan portfolio as it stood on the adoption date, itself the product of
lending decisions made years or decades prior.  The CECL charge is large when
a bank's portfolio contains loans with long remaining maturities, weak
borrower credit quality, or thin collateral coverage, characteristics fixed
by prior underwriting decisions rather than 2023 deposit conditions.  The
cross-sectional regression in Table~\ref{tab:cecl_predictors_cross_section}
confirms that the CECL adjustment is predicted by NPL ratios and historical
allowance levels, with a maximum adjusted $R^2$ of roughly one percent,
and shows no meaningful relationship with deposit structure variables.
This evidence is consistent with treatment intensity being largely independent
of pre-adoption deposit outcomes, though we cannot rule out residual
correlation through channels not captured by the observables in our sample.


\subsection{Information Content, Not Cash Flows}
\label{sec:identification:mechanism}

A central advantage of our design is the separation of information from
concurrent changes in bank financial condition.  Prior depositor-discipline
tests typically exploit realized distress events, rating downgrades, NPL
spikes, or near-failure episodes, in which new information arrives
simultaneously with deterioration in actual cash flows or liquidity
\citep{flannery1996evidence, martinez2001depositors, iyer2016tale}.
In such settings, the response to information is inseparable from the
response to contemporaneous financial deterioration.  The CECL Day-One
adjustment resolves this confound: as discussed in
Section~\ref{sec:background}, it records a pure accounting charge that does
not trigger a cash transfer, alter borrower repayment terms, reduce bank
liquidity, or change the risk of the underlying loans.  A depositor who
withdraws funds or demands a higher rate after observing a large CECL
adjustment is responding to newly disclosed information about credit-risk
composition, not to a concurrent worsening of the bank's actual financial
condition.

The disclosure is also mandatory, quantified, and audited, which distinguishes
it from voluntary disclosure events where the decision to disclose is itself
endogenous to management strategy, and from supervisory rating changes where
the information event is confounded by the behavioral constraints regulators
impose directly.

\subsection{Threats to Identification}
\label{sec:identification:threats}

The principal threat to identification is the coincidence of the primary
adoption quarter (2023Q1) with the failures of Silicon Valley Bank and
Signature Bank in March 2023.  If depositors at high-CECL banks responded
to aggregate stress rather than CECL-specific information, the treatment
effect would be confounded.  Three features of our design limit this concern.
The 2023 stress was driven by unrealized securities losses and duration
mismatches, not by credit risk; Table~\ref{tab:cecl_predictors_cross_section}
shows that unrealized securities losses have no meaningful correlation with
the CECL treatment variable, so high-CECL banks are not disproportionately
exposed to the 2023 stress channel.  Calendar-quarter fixed effects absorb
the aggregate shock common to all community banks in 2023Q1, and the
treatment effect is identified from the differential response at high-CECL
relative to low-CECL banks within that quarter.  Finally, the community banks
in our sample are substantially insulated from the 2023 stress: all have
assets below \$10 billion, carry primarily loan-funded balance sheets, and
lack the concentrated uninsured technology-sector deposit bases that made
SVB uniquely fragile \citep{jiang2024monetary}.  A secondary concern is
differential rate sensitivity: if high-CECL banks were more exposed to the
Federal Reserve's 2022--2023 tightening cycle, their deposit dynamics would
reflect rate pass-through rather than information.  The lagged interest
expense-to-assets control absorbs pre-existing differences in pricing
strategy, and calendar-quarter fixed effects absorb aggregate rate variation.

\subsection{Signal Validation: Market Responses at Large-Bank Adoption}
\label{sec:identification:signal}

A necessary condition for our interpretation is that the CECL adjustment
contained information about bank credit quality that was not already reflected
in prices.  We test this using the cohort of large, publicly traded banks that
adopted CECL in 2020Q1, for which equity prices and CDS spreads are observable
around the adoption date.  We do not use this cohort for
our main tests because the 2020Q1 adoption date coincides with the onset of
the COVID-19 pandemic, which overwhelmed any deposit-market signal with
aggregate liquidity demand and emergency policy interventions.

Figure~\ref{fig:es_large_banks_market_cds} reports event-study estimates of
the differential market response around CECL adoption, using the continuous
\texttt{CECL Adj} ratio as treatment intensity with bank and
calendar-month fixed effects.  Panel~A shows equity market capitalization:
the pre-adoption coefficients across nine lead months are flat and
statistically indistinguishable from zero, ruling out systematic anticipation
by equity investors.  At the adoption month, market capitalization drops
sharply and persistently for banks with larger CECL adjustments.  Panel~B
shows CDS spreads: the pre-adoption pattern is similarly flat, and spreads
widen materially at adoption and remain elevated in subsequent months.
Both markets revised assessments sharply at the moment of disclosure, not
before, consistent with the CECL adjustment conveying genuine information
about credit-risk composition that was not recoverable from prior disclosures.
This evidence supports the interpretation that community bank uninsured
depositors in 2023 were responding to a real, previously unavailable signal,
rather than to a mechanical accounting reclassification.
